\chapter*{Введение} % * не проставляет номер
\addcontentsline{toc}{chapter}{Введение} % вносим в содержание

\textbf{Актуальность темы исследования.} В современных информационных системах ключевую и неотъемлемую роль играют базы данных, обеспечивая надежное хранение, обработку и бесперебойный доступ к данным. От стабильности и производительности системы управления базами данных (СУБД) напрямую зависит корректность выполнения операций, скорость обработки пользовательских запросов и общая доступность информационных ресурсов.

В условиях постоянного роста объема данных и стремительно возрастающей нагрузки на СУБД, особую значимость приобретает задача эффективного мониторинга состояния СУБД.

Надёжная работа СУБД напрямую зависит от своевременного выявления проблем, связанных с производительностью, доступностью компонентов и использованием вычислительных ресурсов. Возникновение подобных отклонений может привести к критическому замедлению работы системы или к полной потере доступа к данным. В этой связи системы мониторинга становятся неотъемлемой частью инфраструктуры, обеспечивая постоянный контроль состояния базы данных и быстрое реагирование на сбои в работе системы. 

Таким образом, актуальность настоящего исследования обусловлена отсутствием бесплатного комплексного решения для мониторинга СУБД Greenplum, обладающего необходимым набором функций для оперативного контроля состояния кластера. Существующие инструменты реализованы в рамках коммерческих продуктов, что ограничивает их доступность, и обуславливает необходимость разработки альтернативного решения.

\textbf{Объектом исследования} является процесс мониторинга технического состояния, производительности и доступности распределённой СУБД Greenplum.

\textbf{Предмет исследования} -- технологии и программные средства автоматизированного мониторинга СУБД Greenplum, обеспечивающие контроль за ключевыми метриками производительности и формирование системы оповещений о критических инцидентах.

\textbf{Цель данной работы} -- разработка комплексной системы мониторинга и оповещений для СУБД Greenplum. 

Для достижения поставленной цели были поставлены следующие \textbf{задачи работы}:
\begin{enumerate}
	\item Исследовать предметную область, изучить существующие решения и определить их достоинства и недостатки;
	\item Выбрать необходимые инструменты и технологии для реализации системы мониторинга;
	\item Спроектировать и реализовать систему мониторинга;
	\item Провести тестирование и апробацию разработанного решения.
\end{enumerate}
